% =====================================================================
% uni-thesis-engine -- 中文模块化毕业论文模板（XeLaTeX / LuaLaTeX）
% 使用 ctex 宏集与 TeX Live 自带的开源 Fandol 字体，无需任何商用字体。
% 所有实验数值、图表与引用均来自 generated/ 下的自动生成文件：
%   generated/auto_metrics.tex     结果宏（\Metric... \Stat... \Diff...）及条件开关
%   generated/auto_tables.tex      booktabs 三线表
%   generated/auto_figures.tex     fig_foundry 矢量图
%   generated/auto_literature.tex  DOI 核验后的相关文献
% 切勿在 chapters/ 中手写实验数值，否则反幻觉审计将拒绝通过。
% =====================================================================
\documentclass[UTF8,zihao=-4,fontset=fandol,oneside]{ctexrep}

\usepackage[margin=2.5cm]{geometry}
\usepackage{amsmath,amssymb}
\usepackage{booktabs}
\usepackage{graphicx}
\usepackage[round,authoryear]{natbib}
\usepackage[hidelinks]{hyperref}

\ctexset{bibname = {参考文献}}

\input{generated/auto_metrics.tex}

\title{\ThesisTitle}
\author{作者姓名\\ \small 指导教师：导师姓名}
\date{\today}

\begin{document}
\maketitle

\begin{abstract}
本文研究课题“\ThesisTitle”。研究构建了一条可复现的自动化流程：检索并以 DOI 核验相关文献，
依据声明式课题配置执行全部实验，按数据假设自动路由统计检验，并通过自动生成的宏将全部结果注入正文。
在主要指标（\MetricPrimaryName）上，本文配置（\MetricPrimaryProposedName）达到
\MetricPrimaryProposedMean{} $\pm$ \MetricPrimaryProposedStd{} \MetricPrimaryUnit，最强对照方案
（\MetricPrimaryRunnerUpName）为 \MetricPrimaryRunnerUpMean（$p$ \StatPrimaryPval，\StatPrimaryTest）。
\end{abstract}

\tableofcontents

\chapter{绪论}
\label{chap:intro}

\section{研究背景与意义}
本文围绕课题“\ThesisTitle”展开。可信的实证结论要求每一个报告数值都能追溯到实际执行的实验，
每一条引用都能对应到可核验的出版物。因此，本研究以一条自动化、可审计的科研流程为核心组织全部工作。

\section{研究目标}
\begin{itemize}
  \item 建立完全由声明式课题配置驱动、可复现的实验方案；
  \item 在\MetricPrimaryName{}及各次要指标上，以满足前提假设的统计检验量化本文配置相对其他方案的表现；
  \item 正文只报告实验生成的数值，并以溯源哈希将论文与结果包相互绑定。
\end{itemize}

\section{论文结构}
第~\ref{chap:related}~章综述经 DOI 核验的相关文献；第~\ref{chap:method}~章介绍研究方法与统计方案；
第~\ref{chap:experiments}~章报告实验结果；第~\ref{chap:conclusion}~章总结全文。

\chapter{相关工作}
\label{chap:related}

以下文献检索自 OpenAlex 与 Crossref，且仅在其 DOI 元数据（题名、作者与年份）经 Crossref 核实后保留。
文献局限性在摘要可得时逐字摘录自原文摘要，否则标记为待人工审阅。

\input{generated/auto_literature.tex}

\chapter{研究方法}
\label{chap:method}

\section{实验方案}
研究的全部设置（数据来源、对比方法、评价指标、重复次数与随机种子）均在课题配置
\texttt{\ThesisProfileId} 中声明。每种配置均进行独立重复试验，每次运行都记录在结果包中，
并附带源代码提交号与输入数据的 SHA-256 摘要。

\ifAutoHasRsm
\section{响应面方法}
采用由析因点、轴向点与重复中心点组成的中心复合设计考察各工艺因子，以二阶多项式
\begin{equation}
  \hat{y} = \beta_{0} + \sum_{i} \beta_{i} x_{i} + \sum_{i} \beta_{ii} x_{i}^{2}
            + \sum_{i<j} \beta_{ij} x_{i} x_{j}
\end{equation}
描述响应，并用普通最小二乘法拟合。模型显著性由 $F$ 检验评估，充分性由相对纯误差的失拟检验评估，
驻点性质则由二次型特征值判定。
\fi

\ifAutoHasAblation
\section{消融设计}
在相同条件下评估完整系统的累加式或逐一移除式配置，以分离各核心模块的贡献。
\fi

\section{统计分析}
对每个指标，先以 Shapiro--Wilk 检验考察正态性、以 Levene 检验考察方差齐性。两组比较依据诊断结果
路由至配对或 Student $t$ 检验、Welch $t$ 检验、Wilcoxon 符号秩检验或 Mann--Whitney $U$ 检验，
并给出 Cohen's $d$ 或秩二列相关效应量；多组比较采用单因素方差分析加 Tukey HSD 事后检验，
或 Kruskal--Wallis $H$ 检验。显著性标记的阈值见各表注。

\chapter{实验与结果分析}
\label{chap:experiments}

\section{概述}
本章报告课题配置 \texttt{\ThesisProfileId} 中声明的全部实验结果。下文所有数值均为由签名结果包
自动生成的宏，表格与图形亦由同一流程产出。

\input{generated/auto_tables.tex}

\ifAutoHasComparison
\section{与基线方法的对比}
表~\ref{tab:auto-comparison}~汇总了各项指标。在\MetricPrimaryName{}上，本文方法
（\MetricPrimaryProposedName）达到 \MetricPrimaryProposedMean{} $\pm$ \MetricPrimaryProposedStd{}
\MetricPrimaryUnit（置信区间 \MetricPrimaryProposedCiLow{} 至 \MetricPrimaryProposedCiHigh），
最强基线（\MetricPrimaryRunnerUpName）为 \MetricPrimaryRunnerUpMean{} $\pm$ \MetricPrimaryRunnerUpStd。
两者相差 \DiffPrimaryAbs（相对提升 \DiffPrimaryRelPct\,\%），经\StatPrimaryTest{}检验
$p$ \StatPrimaryPval（效应量 \StatPrimaryEffectName{} 为 \StatPrimaryEffect；每种方法 \StatPrimaryN{} 次运行）。
全部方法间的\StatPrimaryOmnibusTest{}给出 $p$ \StatPrimaryOmnibusPval。
\fi

\ifAutoHasAblation
\section{消融实验}
表~\ref{tab:auto-ablation}~列出了各消融配置。完整系统在\MetricPrimaryName{}上达到
\MetricPrimaryProposedMean{} $\pm$ \MetricPrimaryProposedStd{} \MetricPrimaryUnit，最优的精简配置
（\MetricPrimaryRunnerUpName）为 \MetricPrimaryRunnerUpMean，提升 \DiffPrimaryAbs（\DiffPrimaryRelPct\,\%；
\StatPrimaryTest，$p$ \StatPrimaryPval）。各配置间的\StatPrimaryOmnibusTest{}给出 $p$ \StatPrimaryOmnibusPval。
\fi

\ifAutoHasRsm
\section{响应面寻优}
各设计试验点列于表~\ref{tab:auto-design}。拟合模型（表~\ref{tab:auto-coef}）为
\begin{equation}
  \StatRsmPrimaryEquation
\end{equation}
其中 \StatRsmPrimaryFactorLegend。表~\ref{tab:auto-anova}~的方差分析给出
$R^{2}$ = \StatRsmPrimaryRsq，调整 $R^{2}$ = \StatRsmPrimaryRsqAdj（$F$ = \StatRsmPrimaryFvalue，
$p$ \StatRsmPrimaryPval），失拟检验 $p$ \StatRsmPrimaryLackOfFitPval。驻点为\StatRsmPrimaryOptNature，
位于 \StatRsmPrimaryOptPoint，预测响应为 \StatRsmPrimaryOptResponse{} \MetricPrimaryUnit。
图~\ref{fig:auto-rsm}~展示了响应面形态。
\fi

\section{图示}
\input{generated/auto_figures.tex}

\chapter{结论与展望}
\label{chap:conclusion}

\section{全文总结}
本文以全程可追溯的流程研究了“\ThesisTitle”。在主要指标（\MetricPrimaryName）上，本文配置
（\MetricPrimaryProposedName）达到 \MetricPrimaryProposedMean{} \MetricPrimaryUnit，较
\MetricPrimaryRunnerUpName{} 提升 \DiffPrimaryAbs（$p$ \StatPrimaryPval）。
\ifAutoHasRsm
响应面分析在 \StatRsmPrimaryOptPoint{} 处找到\StatRsmPrimaryOptNature，预测\MetricPrimaryName{}为
\StatRsmPrimaryOptResponse。
\fi

\section{不足与展望}
本文结论受限于课题配置所声明的数据来源与工作范围。后续工作应在独立数据集上验证结论；
当最优点靠近边界时扩展因子范围；并考察与实际部署相关的更多指标。


\bibliographystyle{plainnat}
\bibliography{refs}

\end{document}
